\subsection{Examples of Springer fibers} 

\subsubsection{} When $e=0$, $\sB_{e}=\sB$.  

\subsubsection{Regular nilpotent elements}\label{sss:regular Spr} The unique dense open $G$-orbit consists of {\em regular nilpotent elements}, i.e., those $e$ such that $\dim C_{G}(e)=r$. When $e$ is a regular nilpotent element, $\sB_{e}$ is a single point: there is a unique Borel subgroup of $G$ whose Lie algebra contains $e$. What is this Borel subgroup? 

By the Jacobson-Morosov theorem, we may extend $e$ to an $\sl_{2}$-triple $(e,h,f)$ in $\frg$. The adjoint action of $h$ on $\frg$ has integer weights, and it decomposes $\frg$ into weight spaces $\frg(n)$, $n\in\ZZ$. Let $\frb=\oplus_{n\geq 0}\frg(n)$. This is a Borel subalgebra of $\frg$, and the corresponding Borel subgroup is the unique point in $\sB_{e}$.

When $e$ is regular,  the fiber $\tB_{e}$ of the Grothendieck alteration is a non-reduced scheme whose underlying reduced scheme is a point. The coordinate ring of $\tB_{e}$ is isomorphic to the coinvariant algebra $\Sym(\frt^{*})/(\Sym(\frt^{*})^{W}_{+})$, which is, interestingly, also isomorphic to the cohomology ring with $k$-coefficients of the complex flag variety $\calB_{\CC}$.

\subsubsection{} When $G=\SL(V)$ for some vector space $V$ of dimension $n$, $\sB$ is the moduli space of full flags $0=V_{0}\subset V_{1}\subset V_{2}\subset\cdots\subset V_{n-1}\subset V_{n}=V$. The Springer fiber $\sB_{e}$ consists of those flags such that $eV_{i}\subset V_{i-1}$.



\subsubsection{}\label{sss:SL3 subreg} Consider the case $G=\SL_{3}$ and $e=\left(\begin{array}{ccc} 0 & 0 & 1\\  0 & 0& 0 \\  0 & 0& 0 \end{array}\right)$ under the standard basis $\{v_{1}, v_{2}, v_{3}\}$ of $V$. Then $\sB_{e}$ is the union of two $\PP^{1}$s: the first $\PP^{1}$ consisting of flags $0\subset V_{1}\subset \jiao{v_{1}, v_{2}}\subset V$ with varying $V_{1}$ inside the fixed plane $\ker(e)=\jiao{v_{1},v_{2}}$; the second $\PP^{1}$ consisting of flags $0\subset \jiao{v_{1}}\subset V_{2}\subset V$ with a varying $V_{2}$ containing $\Im(e)=\jiao{v_{1}}$. 



\subsubsection{} Consider the case $G=\SL_{4}$ and the nilpotent element $e: v_{3}\mapsto v_{1}\mapsto 0, v_{4}\mapsto v_{2}\mapsto 0$ under the standard basis $\{v_{1},v_{2},v_{3},v_{4}\}$. If a flag $0\subset V_{1}\subset V_{2}\subset V_{3}\subset V$ is in $\sB_{e}$, then $V_{1}\subset \ker(e)=\jiao{v_{1},v_{2}}$, and $V_{3}\supset\Im(e)=\jiao{v_{1},v_{2}}$. We denote $H:=\jiao{v_{1},v_{2}}$. There are two cases:
\begin{enumerate}
\item $V_{2}=H$. Then we may choose $V_{1}$ to be any line in $H$ and $V_{3}$ to be any hyperplane containing $V_{2}$. We get a closed subvariety of $\sB_{e}$ isomorphic to $\PP(H)\times \PP(H)\cong \PP^{1}\times \PP^{1}$. We denote this closed subvariety of $\sB_{e}$ by $C_{1}$. 
\item $V_{2}\neq H$. This defines an open subscheme $U$ of $\sB_{e}$. Suppose the line $V_{1}\subset H$ is spanned by  $av_{1}+bv_{2}$ for some $[a:b]\in\PP(H)$, then the image of $V_{2}$ in $V/H=\jiao{v_{3},v_{4}}$ is spanned by $av_{3}+bv_{4}$. Fixing $V_{1}$, the choices of $V_{2}$ are given by $\Hom(\jiao{av_{3}+bv_{4}}, H/V_{1})\cong\Hom(V_{1},H/V_{1})$. Once $V_{2}$ is fixed, $V_{3}=V_{2}+H$ is also fixed. Therefore $U$ is isomorphic to the total space of the line bundle $\calO(2)$ over $\PP(H)\cong \PP^{1}$.
\end{enumerate}
From the above discussion we see that $\sB_{e}$ has dimension $2$, $C_{1}$ is an irreducible component of $\sB_{e}$ and so is the closure of $U$, which we denote by $C_{2}$. We have $C_{1}\cup C_{2}=\sB_{e}$ and $C_{1}\cap C_{2}$ is the diagonal inside $C_{1}\cong\PP^{1}\times\PP^{1}$. 
 
\subsubsection{Components of type $A$ Springer fibers}\label{sss:Spr type A} When $G=\SL_{n}=\SL(V)$, Spaltenstein \cite{Spa76} and Steinberg \cite[]{St} gave a description of the irreducible components of  $\sB_{e}$ using standard Young tableaux of size $n$. This will be relevant to the Springer correspondence that we will discuss later, see \S\ref{sss:A Spr corr}. Below we follow the presentation of \cite[Ch II, \S5]{Spaltenstein}.

Fix a nilpotent element $e\in\sN$ whose Jordan type is a partition $\l$ of $n$. This means, if the partition $\l$ is $n=\l_{1}+\l_{2}+\cdots$, $e$ has Jordan blocks of sizes $\l_{1},\l_{2},\cdots$. We shall construct a (non-algebraic) map $\sB_{e}\to ST(\l)$, where $ST(\l)$ is the discrete set of standard Young tableau for the partition $\l$. For each full flag $0=V_{0}\subset V_{1}\subset\cdots\subset V_{n-1}\subset V_{n}=V$ such that $eV_{i}\subset V_{i-1}$, $e$ induces a nilpotent endomorphism of $V/V_{n-i}$. Let $\mu_{i}$ be the Jordan type of the $e$ on $V/V_{n-i}$, then $\mu_{i}$ is a partition of $i$. The increasing sequence of partitions $\mu_{1},\mu_{2},\cdots, \mu_{n}=\l$ satisfies that $\mu_{i}$ is obtained from $\mu_{i-1}$ by increasing one part of $\mu_{i-1}$ by 1 (including creating a part of size $1$). This gives an increasing sequence $Y_{1},\cdots, Y_{n}=Y(\l)$ of subdiagrams of the Young diagram $Y(\l)$ of $\l$. We label the unique box in $Y_{i}-Y_{i-1}$ by $i$ to get a standard Young tableau.
 
Spaltenstein \cite[Ch II, Prop 5.5]{Spaltenstein} showed that the closure of the preimage of each standard Young tableaux in $\sB_{e}$ is an irreducible component. Moreover, all irreducible components of $\sB_{e}$ arise in this way and they all have the same dimension $d_{e}=\frac{1}{2}\sum_{i}\l^{*}_{i}(\l^{*}_{i}-1)$, where $\l^{*}$ is the conjugate partition of $\l$. In particular, the top dimensional cohomology $\cohog{2d_{e}}{\sB_{e}}$ has dimension equal to $\#ST(\l)$, which is also the dimension of an irreducible representation of the symmetric group $S_{n}$. This statement is a numerical shadow of the Springer correspondence, which says that $\cohog{2d_{e}}{\sB_{e}}$ is naturally an irreducible representation of $S_{n}$. 

Spaltenstein \cite[Ch II, Prop 5.9]{Spaltenstein} also showed that there exists a stratification of $\sB$ into affine spaces such that $\sB_{e}$ is a union of strata. This implies that the restriction map on cohomology $\cohog{*}{\sB}\to \cohog{*}{\sB_{e}}$ is surjective. 


\subsubsection{}\label{sss:Sp4} Consider the case $G=\Sp(V)$ for some symplectic vector space $V$ of dimension $2n$, then $\sB$ is the moduli space of full flags
\begin{equation*}
0=V_{0}\subset V_{1}\subset V_{2}\subset\cdots\subset V_{n}\subset\cdots\subset V_{2n-1}\subset V_{2n}=V
\end{equation*}
such that $V^{\bot}_{i}=V_{2n-i}$ for $i=1,\cdots, n$. The Springer fiber $\sB_{e}$ consists of those flags such that $eV_{i}\subset V_{i-1}$ for all $i$.

Consider the case where $\dim V=4$. We choose a basis $\{v_{1},v_{2},v_{3},v_{4}\}$ for $V$ such that the symplectic form $\om$ on $V$ satisfies $\om(v_{i},v_{5-i})=1$ if $i=1$ and $2$, and $\om(v_{i},v_{j})=0$ for $i+j\neq 5$. Let $e$ be the nilpotent element in $\frg=\sp(V)$ given by $e:v_{4}\mapsto v_{1}\mapsto 0, v_{3}\mapsto 0,v_{2}\mapsto 0$. Then a flag $0\subset V_{1}\subset V_{2}\subset V_{1}^{\bot}\subset V$ in $\sB_{e}$ must satisfy $\jiao{v_{1}}\subset V_{2}\subset\jiao{v_{1},v_{2},v_{3}}$, and this is the only condition for it to lie in $\sB_{e}$ (Exercise \ref{exI:211Sp4}). Such a $V_{2}$ corresponds to a line $V_{2}/\jiao{v_{1}}\subset \jiao{v_{2},v_{3}}$, hence a point in $\PP^{1}=\PP(\jiao{v_{2},v_{3}})$. Over this $\PP^{1}$ we have a tautological rank two bundle $\calV_{2}$ whose fiber at $V_{2}/\jiao{v_{1}}$ is the two-dimensional vector space $V_{2}$. The further choice of $V_{1}$ gives a point in the projectivization of $\calV_{2}$. The exact sequence $0\to\jiao{v_{1}}\to V_{2}\to V_{2}/\jiao{v_{1}}\to0$ gives an exact sequence of vector bundles $0\to\calO\to\calV_{2}\to\calO(-1)\to0$ over $\PP^{1}$. Therefore $\calV_{2}$ is isomorphic to $\calO(-1)\oplus\calO$, and $\sB_{e}\cong \PP(\calO(-1)\oplus\calO)$ is a Hirzebruch surface.



  
\subsubsection{Subregular Springer fibers}\label{sss:subreg} The example considered in \S\ref{sss:SL3 subreg} is a simplest case of a {\em subregular Springer fiber}. There is a unique nilpotent orbit $\calO_{\subreg}$ of codimension 2 in $\sN$, which is called the {\em subregular nilpotent orbit}. For $e\in\calO_{\subreg}$, it is known that $\sB_{e}$ is a union of $\PP^{1}$'s whose configuration we now describe. We may form the dual graph to $\sB_{e}$ whose vertices are the irreducible components of $\sB_{e}$ and two vertices are joined by an edge if the two corresponding components intersect (it turns out that they intersection at a single point).

For simplicity assume $G$ is of adjoint type. Let $G'$ be another adjoint simple group whose type is defined as follows. When $G$ is simply-laced, take $G'=G$. When $G$ is of type $B_{n}$, $C_{n}$, $F_{4}$ and $G_{2}$, take $G'$ to be of type $A_{2n-1}$, $D_{n+1}$, $E_{6}$ and $D_{4}$ respectively. One can show that $\sB_{e}$ is always a union of $\PP^{1}$ whose dual graph is the Dynkin diagram of $G'$. The rule in the non-simply-laced case is that each long simple root corresponds to 2 or 3 $\PP^{1}$'s while each short simple root corresponds to a unique $\PP^{1}$. Such a configuration of $\PP^{1}$'s is called a Dynkin curve, see \cite[\S3.10, Definition and Prop 2]{StConj} and \cite[\S6.3]{Slodowy}.

For example, when $G$ is of type $A_{n}$, then $\sB_{e}$ is a chain consisting of $n$ $\PP^{1}$'s:  $\sB_{e}=C_{1}\cup C_{2}\cup\cdots\cup C_{n}$ with $C_{i}\cap C_{i+1}$ a point and otherwise disjoint.


Brieskorn \cite{Brieskorn}, following suggestions of Grothendieck, related the singularity of the nilpotent cone along the subregular orbits with Kleinian singularities, and he also realized the semi-universal deformation of this singularity inside $\frg$. Assume $G$ is simply-laced. One can construct a transversal slice $S_{e}\subset\sN$ through $e$ of dimension two such that $S_{e}$ consists of regular elements except $e$. Then $S_{e}$ is a normal surface with a Kleinian singularity at $e$ of the same type as the Dynkin diagram of $G$. \footnote{A Kleinian singularity is a surface singularity analytically isomorphic to the singularity at $(0,0)$ of the quotient of $\AA^{2}$ by a finite subgroup of $\SL_{2}$.}.  The preimage $\tilS_{e}:=\pi^{-1}(S_{e})\subset\tN$ turns out to be the minimal resolution of $S_{e}$, and hence $\sB_{e}$ is the union of exceptional divisors. Upon identifying the components of $\sB_{e}$ with simple roots of $G$, the intersection matrix of the exceptional divisors is exactly the negative of the Cartan matrix of $G$. Slodowy \cite{Slodowy} extended the above picture to non-simply-laced groups, and we refer to his book \cite{Slodowy} for a beautiful account of the connection between the subregular orbit and Kleinian singularities. 



